\documentclass[11pt]{report}
\input{../pri-end/T0_preamble_standalone_A4_De}
% T0_preamble_patches.tex — LOKALER PLATZHALTER
% Im Repo durch die offizielle Version ersetzen.


\usepackage{longtable}
\usepackage{booktabs}

\hypersetup{
    pdftitle={251 · FFGFT und Vopsons Computational Universe},
    pdfauthor={Johann Pascher},
    pdfsubject={Vergleich FFGFT mit Infodynamik, Parallelen und Unterschiede}
}

\fancyhead[L]{\small\textcolor{t0gray}{FFGFT · Vopson-Vergleich}}
\fancyhead[R]{\small\textcolor{t0gray}{Johann Pascher · 2026}}

\title{%
    {\LARGE\bfseries\color{t0blue}
     FFGFT und Vopsons Computational Universe}\\[0.5em]
    {\large Parallelen und Unterschiede}\\[1em]
    {\normalsize Johann Pascher \quad $\cdot$ \quad Mai 2026}\\[0.3em]
    {\small \url{https://github.com/jpascher/T0-Time-Mass-Duality}}\\[0.3em]
    {\small Zenodo: \href{https://doi.org/10.5281/zenodo.20117635}{10.5281/zenodo.20117635} (v1.1.0)}
}
\date{Mai 2026}

\begin{document}
    \maketitle
    % ==============================================================================
% 251_FFGFT_Vopson_Vergleich_De_ch.tex
% FFGFT und Vopsons Computational Universe -- Parallelen und Unterschiede
% Autor: Johann Pascher · Mai 2026
% ==============================================================================

\section*{Abstract}

Dieses Dokument vergleicht FFGFT (Fundamentale Fraktale Geometrische Feldtheorie,
T0-Zeit-Masse-Dualität) mit Vopsons Computational-Universe-Hypothese und dem
Zweiten Gesetz der Infodynamik. Beide Ansätze starten von unterschiedlichen
Primitiven --- FFGFT von der 4D-Torusgeometrie~$T^4$, Vopson von
Shannons Informationstheorie und Landauers Prinzip --- und
gelangen zu strukturell sehr ähnlichen Schlussfolgerungen. Die Übereinstimmungen
sind zahlreich und substanziell; der wesentliche Unterschied liegt in der
Richtung der Ableitung und in einem strukturellen Dissens bezüglich der
Möglichkeit einer Außenperspektive auf das System.

\tableofcontents
\newpage

% ------------------------------------------------------------------------------
\section{Einleitung: Zwei Wege zum selben Ziel}

Melvin Vopson entwickelt seit mehreren Jahren die Hypothese eines
\emph{berechnenden Universums} (Computational Universe) und das
\emph{Zweite Gesetz der Infodynamik}: Information ist eine physikalische Größe,
die Informationsentropie bleibt gleich oder nimmt ab, und Gravitation sowie
Symmetrie lassen sich aus informationstheoretischen Minimierungsprinzipien
ableiten \cite{Vopson2022,Vopson2023,Vopson2024}.

FFGFT beginnt an einem anderen Punkt: der Überzeugung, dass zwischen
Quantenmechanik und Relativitätstheorie kein fundamentaler Widerspruch
existieren kann, weil beide zu gut funktionieren. Die Suche nach der
gemeinsamen geometrischen Struktur führte zum 4D-Identifikationstorus~$T^4$
mit dem einzigen dimensionslosen Parameter~$\xipar = 4/3 \times 10^{-4}$ und
der Grundrelation $\tilde{T} \cdot m = 1$.

Wie die folgende Analyse zeigt, sind die Schlussfolgerungen beider Ansätze in
zentralen Punkten strukturell identisch. Das ist kein Zufall: Zwei unabhängige
Wege, die an derselben Stelle ankommen, sprechen dafür, dass dort etwas ist.

% ------------------------------------------------------------------------------
\section{Parallele 1: Das Minimierungsprinzip}

\subsection{Vopsons Zweites Gesetz der Infodynamik}

Vopson formuliert: Die Informationsentropie in einem abgeschlossenen System
bleibt gleich oder nimmt über die Zeit ab. Das steht im Gegensatz zum zweiten
Hauptsatz der Thermodynamik (physikalische Entropie nimmt zu). Empirische
Belege: Virusmutationen verringern die Genomentropie, Atome besetzen Orbitale
mit niedrigster Informationsentropie (Hundsche Regeln) \cite{Vopson2023}.

\subsection{FFGFT: Resonante Stabilität als geometrisches Minimierungsprinzip}

In FFGFT sind nur Windungskonfigurationen mit ganzzahligen Windungszahlen
$(n_\theta, n_\varphi, n_3, n_4)$ stabil --- die geschlossenen Trajektorien
auf~$T^4$ resonieren, die offenen zerfallen. Die $\xipar$-Rekursionsleiter
selektiert stabile Niveaus: Nur Konfigurationen die die geometrische
Resonanzbedingung erfüllen, persistieren.

Das ist ein geometrisches Minimierungsprinzip: Das Universum tendiert zu den
Konfigurationen, die topologisch nachhaltig sind. Die instabilen zerfallen.
Vopson nennt dies Informationsentropie-Minimierung. FFGFT nennt es
resonante Stabilität. Das Ergebnis ist dasselbe.

\begin{erkenntnis}
Vopsons Zweites Gesetz der Infodynamik und FFGFTs Resonanzstabilität auf~$T^4$
beschreiben dieselbe strukturelle Präferenz: Das System tendiert zu den
geometrisch/informationstheoretisch minimalen stabilen Konfigurationen.
\end{erkenntnis}

% ------------------------------------------------------------------------------
\section{Parallele 2: Symmetrie als emergente Sparsamkeit}

\subsection{Vopson}

Symmetrische Strukturen enthalten weniger Informationsgehalt und ``sparen
Rechenleistung''. Das erklärt, warum das Universum Symmetrie gegenüber
Asymmetrie bevorzugt \cite{Vopson2024}.

\subsection{FFGFT}

Symmetrische Konfigurationen auf~$T^4$ entsprechen geschlossenen
Windungstrajektorien --- sie minimieren die geometrische Wirkung und sind
stabile Fixpunkte der $\xipar$-Rekursion. Asymmetrische Konfigurationen
sind entweder metastabil oder instabil.

Beide Schlussfolgerungen: Symmetrie ist keine Koinzidenz, sondern eine
strukturelle Präferenz des zugrundeliegenden Systems.

% ------------------------------------------------------------------------------
\section{Parallele 3: Gravitation als emergentes Phänomen}

\subsection{Vopson}

Vopson leitet Newtons Gravitationsgesetz aus der Infodynamik ab: Objekte
ziehen sich an, um ihren gemeinsamen Informations-Fußabdruck zu minimieren.
Das Universum ist rechnerisch effizienter mit wenigen großen Objekten
als mit Milliarden einzelner Teilchen \cite{Vopson2022}.

\subsection{FFGFT}

FFGFT leitet Gravitation aus $\tilde{T} \cdot m = 1$ ab: Das T-Feld koppelt
Masse an die lokale Zeitstruktur, was Gravitationseffekte ohne separates
Graviton oder geometrisch separat eingeführte Krümmung erzeugt. Das
Gravitationspotential folgt aus der T-Feld-Geometrie auf~$T^4$.

Zwei unabhängige Herleitungen, dasselbe empirische Gesetz. Das ist die Art
von Konvergenz die bedeutsam ist.

% ------------------------------------------------------------------------------
\section{Parallele 4: Ein nicht direkt beobachtbares Feld erklärt Dunkle Materie}

\subsection{Vopson}

Dunkle Materie ist die gespeicherte Information des Universums --- der
eigentliche ``Programmcode'' der Realität. Nicht direkt beobachtbar als
Materie, aber gravitationswirksam und strukturtragend \cite{Vopson2024}.

\subsection{FFGFT}

Das T-Feld in FFGFT: Nicht direkt beobachtbar als Materie, permeiert den
gesamten Raum, beeinflusst Gravitationsdynamik über $\tilde{T} \cdot m = 1$,
trägt topologische Information durch Windungskonfigurationen. FFGFT benötigt
keine separate Dunkle-Materie-Substanz, weil das T-Feld die beobachteten
Effekte bereits erklärt.

Ob man dies ``gespeicherte Information'' (Vopson) oder ``T-Feld aus
$T^4$-Geometrie'' (FFGFT) nennt --- die funktionale Beschreibung ist dieselbe.

\begin{erkenntnis}
Vopsons ``Dunkle Materie als Information'' und FFGFTs T-Feld erfüllen
strukturell dieselbe Rolle: ein nicht-direkt-beobachtbares Feld das
Gravitation und Struktur erklärt.
\end{erkenntnis}

% ------------------------------------------------------------------------------
\section{Parallele 5: Information und Bewusstsein als Erhaltungsgrößen}

\subsection{Vopson}

Durch die Erhaltungssätze der Physik (Energie und Information können nicht
vernichtet werden) folgert Vopson, dass Bewusstsein bzw. die Information eines
Lebewesens nach dem physischen Tod in irgendeiner Form fortbesteht
\cite{Vopson2024}.

\subsection{FFGFT}

Topologische Invarianten (Windungszahlen) sind in FFGFT strukturell erhalten
--- auch in Extremkonfigurationen wie schwarzen Löchern. Die topologische
Grundinformation (``Grundinformation'' im Sinne von Dok.~250) geht nicht
verloren; sie kann für einen äußeren Beobachter unzugänglich werden
(Zugangsproblem), existiert aber geometrisch weiterhin.

Die Schlussfolgerung ist dieselbe: Information verschwindet nicht.

% ------------------------------------------------------------------------------
\section{Parallele 6: Informationskapazität des Universums}

\subsection{Vopsons Bit-Zahl}

Vopson berechnet die Gesamtinformation des beobachtbaren Universums aus der
Anzahl der Elementarteilchen (~$10^{80}$) mal Bits pro Teilchen.
Ergebnis: ~$10^{80}$ bis $10^{90}$ Bits (je nach Papier) \cite{Vopson2023}.
Das sind die tatsächlich \emph{belegten} Bits.

\subsection{FFGFT: Maximale geometrische Kapazität}

In FFGFT gibt es eine natürliche Obergrenze über $L_0$:

\begin{align}
  L_0 &= \xipar \cdot \ell_P \approx 2{,}2 \times 10^{-39}\,\text{m} \\
  N_\text{max}^\text{(3D)} &= \frac{V_\text{Universum}}{L_0^3}
    \approx \frac{4 \times 10^{80}\,\text{m}^3}{(2{,}2 \times 10^{-39})^3\,\text{m}^3}
    \approx 10^{195}\,\text{Bits} \\
  N_\text{max}^\text{(4D)} &= \frac{V_{4D}}{L_0^4}
    \approx 10^{260}\,\text{Bits}
\end{align}

\begin{longtable}{lll}
\toprule
Größe & Wert & Bedeutung \\
\midrule
Vopson (belegte Bits) & $10^{80}$--$10^{90}$ & Tatsächlich belegte Information \\
FFGFT (max. Kapazität 3D) & $\sim 10^{195}$ & Geometrische Obergrenze \\
FFGFT (max. Kapazität 4D) & $\sim 10^{260}$ & T⁴-Kapazität \\
Füllungsrate (3D) & $\sim 10^{-115}$ & Universum ist fast leer \\
\bottomrule
\end{longtable}

Kein Widerspruch: Vopson zählt den \emph{Inhalt}, FFGFT gibt die
\emph{Behältergröße} an. Das Universum nutzt einen verschwindend kleinen
Bruchteil seiner geometrischen Informationskapazität --- konsistent mit
der beobachteten kosmischen Leere.

% ------------------------------------------------------------------------------
\section{Der Ausgangspunkt: Umgekehrte Richtung}

Der wichtigste Unterschied ist nicht das Ergebnis, sondern die Richtung:

\begin{longtable}{>{\bfseries}p{3cm} p{5cm} p{5cm}}
\toprule
 & Vopson & FFGFT \\
\midrule
Primitiv & Shannon-Information (physikalisch) & $T^4$-Geometrie + $\xipar$ \\
Richtung & Information $\rightarrow$ Geometrie emergiert & Geometrie $\rightarrow$ Information ist Projektion \\
Startpunkt & Landauer-Prinzip, Shannon & QM/RT-Vereinigung, Musiktheorie \\
Gravitation & Aus Informationsminimierung & Aus $\tilde{T} \cdot m = 1$ \\
Dunkle Materie & Gespeicherte Information & T-Feld aus $T^4$ \\
\bottomrule
\end{longtable}

Beide landen an strukturell sehr ähnlichen Orten, weil sie eine gemeinsame
tiefere Realität beschreiben --- von verschiedenen Seiten angegangen.

Das entspricht dem Muster anderer konvergierender Ansätze in der
Wissenschaftsgeschichte: Maxwell und Faraday kamen zur Elektrodynamik von
verschiedenen Seiten; Heisenberg und Schrödinger entwickelten äquivalente
Formulierungen der Quantenmechanik. Die Konvergenz selbst ist ein Signal.

% ------------------------------------------------------------------------------
\section{Der Dissens: Gibt es ein Außen?}

\subsection{Vopsons Position}

Wenn man den Code der Realität entdeckt und versteht, könnte man ihn
theoretisch ``hacken'' --- Teleportation, Überwindung der Lichtgeschwindigkeit,
``gottgleiche Kräfte'' \cite{Vopson2024}. Die Simulationshypothese lässt
ein Außen offen: Das Universum könnte auf etwas anderem laufen.

\subsection{FFGFT: Struktureller Ausschluss}

In FFGFT gibt es kein Außen zu~$T^4$. Jeder Versuch, den Code zu ``hacken'',
ist selbst ein Prozess innerhalb von~$T^4$. Das ist keine praktische Limitation
--- es ist eine geometrische. Die Analogie zu Gödels Unvollständigkeitssatz
(Dok.~248, §9.4) ist direkt: Ein System kann seine eigene Vollständigkeit nicht
aus seinen Axiomen beweisen, und es gibt keinen Standpunkt außerhalb des
Systems von dem aus man es vollständig beschreiben könnte.

Fragen wie ``Was liegt vor~$T^4$?'' oder ``Wer hat~$\xipar$ gewählt?'' sind
in FFGFT nicht wohlgeformt --- nicht weil man sie noch nicht beantwortet
hätte, sondern weil sie eine Außenperspektive voraussetzen die strukturell
nicht existiert (vgl.\ Dok.~248, §9).

\begin{wichtig}
Die Simulationshypothese setzt voraus, dass ein ``Außen'' existiert von dem
aus das Universum simuliert wird. FFGFT schließt das strukturell aus: $T^4$
ist nicht in einem größeren Raum eingebettet. ``Was liegt außerhalb?'' ist
keine beantwortbare Frage --- analog zu ``Was liegt nördlich des Nordpols?''
\end{wichtig}

% ------------------------------------------------------------------------------
\section{Tori-Resonanzen als Bewusstseins-Selektionsmechanismus}

Ein weiterer Berührungspunkt: Vopson sieht das Universum als
informationsverarbeitendes System; Bewusstsein ist ein Teil davon. In FFGFT
sind Tori-Resonanzen bereits ein eingebauter Selektionsmechanismus:

\begin{itemize}
  \item Nur ganzzahlige Windungszahlen sind stabil --- resonante
    Konfigurationen persistieren, nicht-resonante zerfallen.
  \item Das ist ein geometrischer Selektionsmechanismus der auf allen Skalen
    operiert --- von Elementarteilchen bis zu komplexen Systemen.
  \item Bewusstsein emergiert in FFGFT-Perspektive aus der rekursiven
    Rückkopplungsschleife (Sensor → Gedächtnis → Bewegungskorrelation)
    die auf derselben topologischen Basis aufbaut.
\end{itemize}

Rudimentäre Formen dieser Selbstbezüglichkeit existieren bereits in einfachen
T\textsuperscript{4}-Konfigurationen --- analog zu Vopsons Sicht dass
Information in jeder physikalischen Struktur inhärent ist.

% ------------------------------------------------------------------------------
\section*{Zusammenfassung}

\begin{longtable}{>{\bfseries}p{5cm} p{4.5cm} p{4.5cm}}
\toprule
Aspekt & Vopson & FFGFT \\
\midrule
Primitiv & Information (physikalisch) & $T^4$-Geometrie + $\xipar$ \\
Minimierungsprinzip & 2. Gesetz der Infodynamik & $\xipar$-Resonanzstabilität \\
Symmetrie & Informationsminimal & Geometrisch minimal \\
Gravitation & Aus Informationsminimierung & Aus $\tilde{T} \cdot m = 1$ \\
Dunkle Materie & Gespeicherte Information & T-Feld \\
Informationserhaltung & Ja & Ja (topologische Invarianten) \\
Bewusstsein & Information bleibt erhalten & Topologische Ladung erhalten \\
Bits im Universum & $10^{80}$--$10^{90}$ (belegt) & $10^{195}$ (max.\ Kapazität) \\
Außenperspektive & Offen (Simulation möglich) & Strukturell ausgeschlossen \\
Hackbarkeit & Theoretisch möglich & Nein --- kein Außen \\
\bottomrule
\end{longtable}


    \begin{thebibliography}{99}
    \bibitem{Vopson2022}
    Vopson, M. M. (2022). The mass-energy-information equivalence principle.
    \emph{AIP Advances}, 9, 095206.

    \bibitem{Vopson2023}
    Vopson, M. M. (2023). Second law of infodynamics and its implications for the
    simulated universe hypothesis. \emph{AIP Advances}, 13, 105308.

    \bibitem{Vopson2024}
    Vopson, M. M. (2024). Experimental protocol to test the second law of
    infodynamics using a scanning tunnelling microscope. \emph{AIP Advances}, 14, 035205.

    \bibitem{Pascher2026}
    Pascher, J. (2026). \emph{FFGFT v1.1.0}. Zenodo.
    \url{https://doi.org/10.5281/zenodo.20117635}

    \bibitem{Pascher248}
    Pascher, J. (2026). \emph{FFGFT --- Epistemologische Position} (Dok.~248).
    Zenodo. \url{https://doi.org/10.5281/zenodo.20117635}

    \bibitem{Pascher250}
    Pascher, J. (2026). \emph{FFGFT --- Information und schwarze Löcher} (Dok.~250).
    Zenodo. \url{https://doi.org/10.5281/zenodo.20117635}
    \end{thebibliography}
\end{document}
